% [VERIFY] conclusion -- problem.md Key Insight, headline claims C01-C03,C08, constraints.md limitations,
%   trace dead-ends N04 (vanishing-gradient ruled-OUT-as-cause stays hedged), N05 (3x iters), N10 (single-layer F).
% conclusion-limitation: Conclusion = single para, no numbers/citations. Limitations honest + specific.
\section{Conclusion and Limitations}
\label{sec:conclusion}

\noindent\textbf{Conclusion.}
We have presented a residual learning framework that reformulates the layers of
a deep network as learning residual functions with reference to the layer
inputs, realized by parameter-free identity shortcut connections. This
reformulation eases the optimization of very deep networks: whereas plain
networks of increasing depth suffer a degradation in training accuracy, their
residual counterparts are readily optimized and gain accuracy from considerably
increased depth. The framework requires no change to the optimizer and adds no
parameters relative to the plain counterpart, yet it scales to depths far beyond
what was previously trainable, and the depth of the resulting representations
proves to be of central importance not only for image classification but also
for downstream localization-sensitive recognition tasks.

\noindent\textbf{Limitations.}
Several boundaries of the present study should be stated plainly. First, the
claim that the degradation of plain networks stems from optimization difficulty
rather than vanishing gradients is argued from batch-normalized signal variances
and healthy gradient norms, and is supported by the observation that training
longer does not close the gap; it is \emph{not} a formal proof, and the
conjecture of exponentially low convergence rates for plain nets remains open.
Likewise, the premise that stacked nonlinear layers can asymptotically
approximate the residual function is an open question, not an established fact.
Second, the per-iteration training trajectories that motivate the ``throughout
training'' phrasing are reported in the original analysis but are not tabulated
as per-iteration data in the artifact underlying this manuscript; the figures
here therefore present reported endpoint errors rather than redrawn curves.
Third, residual learning eases optimization but does not by itself prevent
overfitting: at extreme depth on a small dataset, the 1202-layer CIFAR network
trains without optimization difficulty yet generalizes worse than a much smaller
residual net, and we did not apply strong regularization to counter this. The
study scales depth only --- it does not investigate width --- and uses
first-order SGD throughout, with no comparison against second-order or
specialized solvers.

\noindent\textbf{Future work.}
Two concrete directions follow. Combining extremely deep residual nets with
explicit regularization such as maxout or dropout would test whether the
overfitting seen at 1202 layers on CIFAR-10 can be removed while retaining the
clean optimization behavior. Pairing residual learning with stronger
regularizers, or with the box-refinement, global-context, and multi-scale
testing extensions used in the competition entries, is a natural path to
quantify how much of the detection gain can be pushed further on top of the
backbone improvement.
